% !TeX root = ../navier-stokes-manuscript.tex
% ============================================================
% 2.4 The Spatial Column Inside Critical Height
% ============================================================

\subsection{The Spatial Column Inside Critical Height}\label{sub:TheSpatialColumnInsideCriticalHeight}

To see exactly how one time derivative reaches into space, take a scalar energy from one spatial level of the mixed jet.
For \(s\ge0\), set
\[
  E_s(t)
  :=
  \frac12\norm{\Lambda^su(t)}_2^2,
\]
so that
\[
  H(t)=E_{1/2}(t).
\]
Keeping \(s\) free before returning to \(1/2\) exposes the mechanism that repeats at every height.

Apply \(\Lambda^s\) to the pressure-complete momentum equation and pair it in \(L^2\) with \(\Lambda^su\).
The temporal term gives
\[
  \pair{\Lambda^su}{\Lambda^s\partial_tu}
  =
  \frac{d}{dt}\frac12\norm{\Lambda^su}_2^2
  =
  E_s'.
\]
Rapid decay and commutation with spatial derivatives give
\[
  \nu\pair{\Lambda^su}{\Lambda^s\Delta u}
  =
  -\nu\norm{\Lambda^{s+1}u}_2^2
  =
  -2\nu E_{s+1}.
\]
The pressure pairing vanishes only after incompressibility and integration by parts have done their work:
\[
  \pair{\Lambda^su}{\Lambda^s\nabla p}
  =
  -\pair{\nabla\cdot\Lambda^su}{\Lambda^sp}
  =
  0.
\]
This global cancellation does not remove pressure from the local equation or from the mixed jet that generated the balance.
Keep the signed nonlinear transfer in its pressure-complete form,
\[
  \mathcal P_s(t)
  :=
  -\pair{\Lambda^su}
  {\Lambda^s\bigl((u\cdot\nabla)u+\nabla p\bigr)}.
\]
The exact balance is
\[
  E_s'
  =
  \mathcal P_s-2\nu E_{s+1}.
\]
Thus one time derivative of \(E_s\) introduces the next spatial energy \(E_{s+1}\) together with the complete signed transfer at height \(s\).

At critical height, the first row reads
\[
  H'
  =
  \mathcal P_{1/2}-2\nu E_{3/2}.
\]
This is already the first spatial ascent, from the scale-invariant \(E_{1/2}\) to \(E_{3/2}\).
It is not a bound for \(E_{3/2}\), because the transfer and dissipative terms can cancel inside the complete value of \(H'\).
The whole row has to be differentiated before its next spatial consequence can be seen.

The balance one step higher is
\[
  E_{3/2}'
  =
  \mathcal P_{3/2}-2\nu E_{5/2}.
\]
Substituting it into the derivative of the first critical row gives
\begin{align*}
  H''
  &=
  \mathcal P_{1/2}'-2\nu E_{3/2}'
  \\
  &=
  \mathcal P_{1/2}'
  -2\nu\mathcal P_{3/2}
  +(2\nu)^2E_{5/2}.
\end{align*}
The second derivative now contains \(E_{5/2}\) and the transfers from both heights \(1/2\) and \(3/2\).
Differentiating once more uses
\[
  E_{5/2}'
  =
  \mathcal P_{5/2}-2\nu E_{7/2}
\]
and yields
\[
  H'''
  =
  \mathcal P_{1/2}''
  -2\nu\mathcal P_{3/2}'
  +(2\nu)^2\mathcal P_{5/2}
  -(2\nu)^3E_{7/2}.
\]
One more response gives
\[
  H^{(4)}
  =
  \mathcal P_{1/2}'''
  -2\nu\mathcal P_{3/2}''
  +(2\nu)^2\mathcal P_{5/2}'
  -(2\nu)^3\mathcal P_{7/2}
  +(2\nu)^4E_{9/2}.
\]
After each time differentiation, the highest spatial energy is one step higher and every earlier transfer is still present.

The pattern is triangular rather than diagonal.
For every integer \(n\ge1\),
\[
  H^{(n)}
  =
  (-2\nu)^nE_{n+1/2}
  +
  \sum_{j=0}^{n-1}
  (-2\nu)^j
  \partial_t^{\,n-1-j}\mathcal P_{j+1/2}.
\]
To pass from \(n\) to \(n+1\), differentiate the displayed identity and substitute
\[
  E_{n+1/2}'
  =
  \mathcal P_{n+1/2}-2\nu E_{n+3/2}.
\]
The substitution produces the next transfer and \(E_{n+3/2}\) with the required coefficient, which proves the identity by induction.
The identity still supplies no estimate by itself.
Even if \(H^{(n)}\) is known, \(E_{n+1/2}\) cannot be isolated without controlling the entire signed transfer sum.
Its appearance in the formula does not bound it.

The homogeneous scale made the critical dilation exact, while continuation is usually read in the inhomogeneous scale.
For \(s>0\),
\[
  \norm{u(t)}_{H^s}^2
  \simeq
  \norm{u(t)}_2^2+\norm{\Lambda^su(t)}_2^2
  =
  \norm{u(t)}_2^2+2E_s(t).
\]
The unforced energy identity supplies the missing low-frequency part:
\[
  \frac12\norm{u(t)}_2^2
  +
  \nu\int_0^t\norm{\nabla u(\tau)}_2^2\,d\tau
  =
  \frac12\norm{u_0}_2^2.
\]
Thus control of \(E_s(t)\) supplies the high-frequency part of an \(H^s\) bound, and an \(L^1\)-in-time bound for \(E_s\) supplies the high-frequency part of an \(L^2_tH^s_x\) bound on each finite interval.

The calculation begins at \(1/2\) because that height is fixed by scaling, and its first derivative reaches \(3/2\).
At the endpoint, the relation runs in the other direction: the whole-terminal Sobolev rows first control the \(1/2\) and \(3/2\) levels, and those levels then control \(H\) and its first dissipative height.

Any finite number of time derivatives reaches only a finite spatial height.
To retain every finite height required by the full derivative family, the same velocity must belong to every finite temporal and spatial Sobolev row.
